\section{Introduction}

Periodic configurations convert an infinite symbolic constraint into a finite
linear problem.  For shifts over an abelian lattice, Fourier characters then
diagonalize convolution, and nullity becomes a count of torsion points on the
zero set of the symbol.  This positive-characteristic viewpoint is developed
systematically by \citet{Zaidenberg2008}.  The discrete Heisenberg group is a
basic nonabelian nilpotent setting in which that character calculation is
incomplete: most of the finite regular representation lies in nonlinear
blocks.

For unit coefficients, the character resultant is Wendt's determinant, whose
relation to finite-field points on Fermat curves is classical
\citep{FordJha1993}.  We therefore do not present the cyclotomic gcd itself as
a new number-theoretic invariant.  Its role is to isolate the abelian stratum
before the nonlinear Heisenberg jump is restored.

We study the finite-field group shift defined by the weighted rule
\begin{equation}\label{eq:local-rule-intro}
 \alpha x_g+\beta x_{ga}+\gamma x_{gb}=0.
\end{equation}
The underlying Heisenberg principal-action framework is established.  In
particular, \citet{GollSchmidtVerbitskiy2014} develop expansiveness and
homoclinic methods for principal actions, and the survey of
\citet{LindSchmidt2015} treats the exact integer element $1+a+b$ as a mixing
example.  Our object is its zero-dimensional coefficient reduction, not a
claim that the group, action framework, or three-term element is new.  The
question here is narrower and arithmetic: what is the exact kernel dimension
after imposing a prime congruence period?

Finite-Heisenberg convolution itself also has a direct computational
literature.  Deundyak and Leonov define left and right convolution on finite
Heisenberg groups, construct the corresponding noncommutative Fourier transform,
and solve convolution equations through the irreducible blocks
\citep{DeundyakLeonov2016}.  That algorithmic framework and its representation
ledger are prior.  It does not calculate the cross-characteristic singular
kernel dimensions of the weighted family in \eqref{eq:local-rule-intro}.

The answer has two qualitatively different pieces.  The one-dimensional
representations behave as in abelian harmonic analysis and contribute a
polynomial gcd degree.  The degree-$\ell$ representations see the clock--shift
matrix
\[
 \alpha I+\beta U+\gamma V.
\]
Its determinant is independent of the nontrivial central character and equals
$\alpha^\ell+\beta^\ell+\gamma^\ell$.  Thus all nonlinear types become
singular together.  Each singular block has nullity one, but its regular
multiplicity is $\ell$; summing over the $\ell-1$ nontrivial central
characters creates a jump of size $\ell(\ell-1)$.

Against the component-by-component owner comparison in
\cref{tab:owner-comparison}, the paper makes three bounded contributions.
\begin{enumerate}
\item For the stated cross-characteristic family, we give an exact
      fixed-dimension formula for every nondegenerate weight
      triple and every pair of distinct primes $p,\ell$, with $\ell$ odd.
\item We separate the character contribution from a nonlinear Fermat-locus
      jump and prove the latter with an exact corank-one statement, not only a
      determinant test.
\item As finite regression evidence, we audit right versus left convolution
      and verify the formula against full $\ell^3$-dimensional matrices in ten
      independent cases.
\end{enumerate}

Finite Heisenberg representation theory, including the Stone--von Neumann
description behind the nonlinear blocks, is standard; we use the formulation
of \citet{GurevichHadani2010} and the direct finite-group classification of
\citet{GrassbergerHormann2001}.  The progress here is the resulting symbolic
fixed-space phase diagram for the weighted cross-characteristic family.  No
priority language is intended: the literature search supporting this internal
draft is bounded.
